\subsubsection{Hlavní skript a jeho role v systému}

Třída \texttt{PlayerController} představuje hlavní skript propojující všechny systémy pohybu hráče. Dědí z \texttt{MonoBehaviour} a přidává se na herní objekt hráče jako komponenta. Tento skript funguje jako centrální uzel, který koordinuje vstupy, fyziku, gravitaci a propojení s dalšími systémy jako dash, střelba a wall jump.

V metodě \texttt{Awake()} probíhá inicializace skriptu. Získávají se reference na komponentu \texttt{Rigidbody2D}, která je esenciální pro fyzikální pohyb, a na ostatní skripty jako \texttt{PlayerDash}, \texttt{PlayerStripeShot} a \texttt{PlayerWall}. Důležitým krokem je nastavení \newline \texttt{gravityScale} na hodnotu 0 na komponentě \texttt{Rigidbody2D}. To je klíčové pro fungování vlastního gravitačního systému, protože Unity vestavěná gravitace je kompletně deaktivována a veškerá gravitace je řešena vlastním kódem v metodě \texttt{ApplyCustomGravity()}.

\begin{lstlisting}[language=csharp]
private void Awake()
{
    rb = GetComponent<Rigidbody2D>();
    rb.gravityScale = 0;
    isFacingRight = true;
    pd = GetComponent<PlayerDash>();
    pss = GetComponent<PlayerStripeShot>();
    pw = GetComponent<PlayerWall>();
}
\end{lstlisting}
